\section{Tensor Omics Module Descriptors (TOXMD)}\label{sec:module_descriptors}

    \subsection{Concept of Modules in Semantic Vector Spaces}

        In Tensor Omics, a \textit{module} is defined as a set of vectors
        \[
            \mathcal{M} = \{ \vec{v}_1, \vec{v}_2, \ldots, \vec{v}_m \}, \quad
            \vec{v}_i \in \mathbb{R}^n,
        \]
        selected according to biological, semantic, or geometric criteria
        (e.g.\ shared molecular function, evolutionary relationship, collinearity,
        local density, or signal strength).

        Modules may originate from vector fields with explicit coordinates
        (e.g.\ spatial or temporal data), but in the most general Tensor Omics setting
        they consist solely of vectors embedded in a semantic omics space.
        Such a module can be represented compactly by the matrix
        \[
            M = [\, \vec{v}_1, \vec{v}_2, \ldots, \vec{v}_m \,] \in \mathbb{R}^{n \times m}.
        \]

        The matrix
        \[
            C = M M^{\mathsf{T}}
        \]
        captures the cumulative second-order structure of the module and serves as the
        basis for several global geometric descriptors.

    \subsection{Canonical Descriptors of a Module}\label{ssec:canonical_module_descriptors}

        From the module matrix \(M\), we define the following canonical descriptors.

        \begin{enumerate}

            \item \textbf{Signal Strength (Energy)}  
                  \[
                      E = \|M\|_{F}
                      = \sqrt{\sum_{i=1}^{n}\sum_{j=1}^{m} M_{ij}^{2}}.
                  \]
                  The Frobenius norm measures the total magnitude of all vectors in the module.
                  In gene expression analysis, this reflects cumulative expression divergence or
                  activity across all genes in the module.

                  To enable comparison across modules of different sizes \(m\), we define the
                  Relative Frobenius Index
                  \[
                      \mathrm{RFI} = \frac{\|M\|_F}{\sqrt{m}},
                  \]
                  which measures average vector magnitude independent of module cardinality.

            \item \textbf{Rank (Subspace Dimensionality)}  
                  \[
                      R = \mathrm{rank}(M).
                  \]
                  The rank quantifies the dimensionality of the subspace spanned by the module.
                  Low rank indicates near-collinearity (coherent behavior), whereas higher rank
                  indicates multi-directional or heterogeneous structure.

            \item \textbf{Principal Direction of Spread}  
                  Let
                  \[
                      C = U \Sigma U^{\mathsf{T}}
                  \]
                  be the eigendecomposition of \(C\).
                  The first eigenvector
                  \[
                      \vec{f} = U_{\cdot 1}
                  \]
                  defines the dominant direction of variance within the module.
                  In semantic omics spaces, this corresponds to the primary axis of coordinated
                  expression or functional divergence.

            \item \textbf{Net Vector (Cumulative Direction)}  
                  \[
                      \vec{s}_{\mathrm{net}} = \sum_{i=1}^{m} \vec{v}_i.
                  \]
                  The net vector summarizes the aggregate directional tendency of the module.
                  Unlike the principal direction, which reflects variance structure, the net
                  vector captures directional bias and asymmetry.
                  Its magnitude vanishes for directionally balanced modules and increases when
                  vectors align coherently in one direction.

            \item \textbf{Angular Dispersion}  
                All vectors are first normalized to unit length,
                \[
                    \hat{\vec{v}}_i = \frac{\vec{v}_i}{\|\vec{v}_i\|}.
                \]
                The spherical mean direction is
                \[
                    \vec{\mu} =
                    \frac{1}{m} \sum_{i=1}^{m} \hat{\vec{v}}_i,
                    \quad
                    R_{\mathcal{M}} = \|\vec{\mu}\|.
                \]
                The scalar \(R_{\mathcal{M}} \in [0,1]\) measures directional coherence.
                Values near 1 indicate strong alignment; values near 0 indicate isotropy.

                A proxy for angular dispersion is given by
                \[
                    \sigma_{\mathrm{ang}} = \sqrt{-2 \ln(R_{\mathcal{M}})},
                \]
                which is monotonic in dispersion and widely used in directional statistics.

            \item \textbf{Axis Footprint (Semantic Support)}  
                  For each axis \(a \in \{1,\ldots,n\}\), let \(v_{a,i}\) denote the \(a\)-th
                  coordinate of vector \(\vec{v}_i\).
                  The footprint of the module along axis \(a\) is defined as
                  \[
                      \left[
                      Q_{\alpha}(v_{a,\cdot}),
                      \;
                      Q_{1-\alpha}(v_{a,\cdot})
                      \right],
                  \]
                  where \(Q_p\) denotes the empirical \(p\)-quantile and
                  \(\alpha \in [0,0.5]\) is a user-defined offset.
                  Setting \(\alpha = 0\) yields the full min-max range; typical values such as
                  \(\alpha = 0.05\) provide robust footprints resistant to outliers.

                  This descriptor localizes the module within the semantic axes of the ambient
                  omics space.

            \item \textbf{Axis Mass (Histogram Representation)}  
                  For each axis, the distribution of \(v_{a,i}\) is summarized by a histogram with
                  a user-defined number of bins \(B\) (default \(B=5\)).
                  Bin boundaries are computed globally from the full ambient space to ensure
                  comparability across modules.

                  Optionally, histogram counts may be weighted by vector magnitudes
                  \( \|\vec{v}_i\| \).
                  Magnitude weighting emphasizes high-expression or high-shift vectors but is
                  generally discouraged in Tensor Omics, as expression magnitude already defines
                  module location in semantic space and additional weighting can obscure
                  directional structure.
                  Unweighted histograms therefore represent the preferred default.

        \end{enumerate}

    \subsection*{Remarks}

        All descriptors are defined for non-zero vectors.
        Modules failing basic validity checks (e.g.\ vanishing angular coherence or
        degenerate dimensionality) should return informative error codes rather than
        numerical values.
        Together, these descriptors form a multi-scale geometric and semantic summary of
        modules in Tensor Omics.

    \subsection{Vector Sets versus Vector Fields}

        A \textit{vector set} (e.g.\ all expression vectors of a gene family) consists 
        only of vectors and their distribution in space.  

        When vectors additionally carry explicit coordinates---as in shift vectors
        originating at gene-family centroids---they form a \textit{vector field}.  In
        such cases, Mod-4-Desc can be applied either to the coordinates or to the
        vectors themselves, yielding complementary information:

    \begin{itemize}
        \item \textbf{On the coordinates} $\vec{x}_i$:  
        the descriptors characterize the geometry and spatial extent of the region 
        in Tensor Omics space, answering questions like:
        \begin{itemize}
            \item Which semantic axes dominate the spatial extent of the field?
            \item Is the field elongated or isotropic (rank)?
            \item Does the region itself curve or spiral through Tensor space (rotation)?
        \end{itemize}

        \item \textbf{On the vectors} $\vec{v}_i$:  
        the descriptors characterize the dynamics or changes that occur within that region:
        \begin{itemize}
            \item Which semantic axes dominate expression changes?
            \item How strong are those changes (energy)?
            \item Are changes co-directed or diverse (rank)?
            \item Do they form coherent rotational flows or transitions (rotation)?
        \end{itemize}

    \end{itemize}

    \textbf{Examples:}  
    In transcriptomics, coordinates correspond to gene-family positions in 
    expression space, while shift vectors represent changes such as stress 
    responses.  
    In physics, coordinates represent spatial positions, and vectors represent wind 
    velocities or magnetic field directions.


    \subsection{Tracing a Module along its Semantic Axis}

        Assume a module has an associated \textit{semantic axis}
        \(\vec{d}_{\text{sem}}\) (see Section~\ref{sec:semantic_axis} for its definition), 
        obtained for example by the Mjölnir manifold learning method.  
        We can then analyze how Mod-4-Desc quantities evolve along this axis.

        A Gaussian kernel \(K(\vec{x}; \sigma)\) is moved along \(\vec{d}_{\text{sem}}\).  
        At each kernel center, the local subset of vectors 
        \( \{\vec{v}_i\} \) within the kernel’s support is collected into a matrix \(M\), 
        and the four TOX Field Descriptors are computed locally:
        \[
            \{E(\vec{x}), R(\vec{x}), \vec{f}(\vec{x}), \omega(\vec{x})\}.
        \]

        Plotting these quantities along the semantic axis reveals how 
        field energy, diversity, and rotational behavior change over semantic space 
        (e.g., developmental time, aging trajectory, or stress gradient).  
        This procedure generalizes naturally to 2D semantic surfaces, 
        where the descriptors can be visualized as tensor fields over the manifold.

    \subsection{Interpretation and Scientific Questions}

        These descriptors allow a model-free geometric interpretation of high-dimensional data:
        \begin{itemize}
            \item \textbf{Energy:} How strong are the local signals or shifts?
            \item \textbf{Rank:} How many independent biological factors are active?
            \item \textbf{Flow:} In which semantic direction does the process move?
            \item \textbf{Rotation:} Are changes linear, oscillatory, or cyclic?
        \end{itemize}

        Together, they enable systematic comparison of complex omics or physical fields, 
        revealing patterns of coherence, divergence, and functional rotation 
        in spaces where traditional differential or statistical models fail.